\section{Related Work}
\label{sec:related}

We review the two categories of Section~\ref{sec:intro} in turn, and then the scheduling theory
CLASP builds on.

\subsection{Physical Link Design for Multi-User Media Delivery}
\label{sec:related-link}

In a cellular downlink, an image is compressed once by a standard codec such as High Efficiency
Video Coding (HEVC) \cite{Sullivan12}, and its bits are carried by a physical layer that adapts to
the channel. Adaptive modulation selects the constellation and the transmit power from the channel state, and
its bit error rate depends on the error and the delay of the channel estimate it adapts to
\cite{Goldsmith97}; in Long Term Evolution (LTE), the receiver reports the channel as a 4-bit CQI
\cite{TS36213}. HARQ retransmits a failed block, and the receiver either combines the noisy copies
\cite{Chase85} or accumulates incremental redundancy; the throughput of such protocols over a
fading Gaussian collision channel has closed forms \cite{Caire01}. Across users, a scheduler
exploits multiuser diversity by serving each user when its instantaneous channel quality is near
its peak, as PF scheduling does \cite{Viswanath02}.

To evaluate such systems without simulating the decoding of every block, system-level simulators
predict a block's
error rate from the signal-to-interference-plus-noise ratios (SINRs) of all the resources it
occupies. An effective SINR mapping compresses them into one effective SNR, from which the block
error rate (BLER) is read off the link-level curve of an additive white Gaussian noise channel.
EESM was derived for orthogonal frequency-division multiplexing (OFDM) from the union--Chernoff
bound, with one parameter tuned per modulation and coding scheme \cite{Ericsson03}; MIESM maps the
SINRs through the mutual information per received symbol or per bit \cite{EMD09}; and the
5G New Radio abstraction of the ns-3 simulator calibrates EESM and obtains SINR--BLER tables with
a link-level simulator \cite{Lagen20}. Both mappings predict whether a block decodes;
Section~\ref{sec:eval-abstraction} tests them as predictors of a learned code's distortion.

Media-aware schedulers value data by the distortion it removes. RaDiO reduces the rate-distortion
optimized streaming of a presentation to the error-cost optimized transmission of each data unit,
solved as a Markov decision process over the unit's transmission opportunities with a Lagrange
multiplier pricing the rate, and essentially achieves the operational distortion-rate function of
the source at the capacity of the packet erasure channel \cite{Chou06}. In a wireless downlink with
several users, gradient-based schedulers set each user's rate from its channel quality and the
gradient of a utility built from the distortion of its received video \cite{Pahalawatta07}.
Hosseini and Panwar model scalable video streaming to many users over Markov channels as a
restless bandit, relax it to a linear program whose value they plot beside their policies, and
improve network utilization by up to a factor of two over PF \cite{Hosseini19}. For packets that
expire at the end of each frame, Hou, Borkar and Kumar characterize the feasible
per-client delivery ratios and show that serving the client with the largest delivery debt first
fulfills every feasible set \cite{HouBorkar09}. All of these treat a block or a packet as decoded
or lost, so the channel affects quality only through how many of them get through, the source of
the cliff effect that learned codes avoid \cite{Bourtsoulatze19}. Our link abstraction takes the
place of an effective SNR mapping for a learned code: it predicts distortion rather than a
probability of decoding, and so it can value every slot a user receives.

\subsection{JSCC Design}
\label{sec:related-jscc}

Deep JSCC trains a convolutional encoder and decoder through a non-trainable channel layer. It
outperforms JPEG or JPEG2000 compression with a capacity-achieving channel code at low SNR and
small bandwidth, has no cliff effect, and degrades gracefully when the SNR departs from its
training value \cite{Bourtsoulatze19}. Later codes adapt one model to the channel and the budget.
Channel-wise attention scales the features by the SNR, so that one model serves a range of SNRs
\cite{Xu22}; nonlinear transform source-channel coding learns an entropy model of its latent
representation and adapts the transmission rate to it \cite{Dai22}; Swin Transformer backbones
adapt to the channel state and the target rate within a single model \cite{Yang25}; and
predictive and adaptive deep coding (PADC) predicts each image's PSNR from its content, the SNR
and the compression ratio, and sends it at the least compression ratio that meets a quality
constraint \cite{Zhang23}. Other codes exploit the structure of the channel: OFDM over multipath
fading \cite{YangBianKim22}, multiple-input multiple-output (MIMO) channels, with robustness to
channel estimation errors \cite{Wu24m}, and channel output feedback, which refines the receiver's
estimate block by block and shortens variable-length transmission \cite{Kurka20,Wu24f}.
DeepJSCC-$l$ sends an image in layers, as successive refinements or as multiple descriptions,
with negligible loss against a single transmission \cite{Kurka21}. A tutorial on semantic and
task-oriented communications is given in \cite{Gunduz23}.

These codes are compared with separation-based transmission on single links: against JPEG or JPEG2000 with a
capacity-achieving code \cite{Bourtsoulatze19}, against Better Portable Graphics (BPG) with
low-density parity-check (LDPC) codes over OFDM, with 2.5--4~dB of SNR gain at equal quality
\cite{YangBianKim22}, and against BPG with 5G LDPC codes \cite{Yang25}. In some of these
comparisons the digital side is idealized, with capacity-achieving codes and exact channel
knowledge at the receiver \cite{Bourtsoulatze19,Kurka20}; Section~\ref{sec:eval-csi} compares
instead with CQI-driven link adaptation and HARQ under erroneous and aged reports.

Linear analog transmission is the precedent for our link abstraction. SoftCast makes the network stack
act as a linear transform of the pixels, so that channel noise reduces sharpness instead of
causing glitches, and one multicast stream gives each receiver a quality commensurate with its
channel \cite{Jakubczak11}. For a linear code, the distortion is a sum of per-component MMSE
terms by construction; the learned code of Section~\ref{sec:method-abstraction} acquires this
form, but only when it is trained for slot-varying SNR.

To serve many users, semantic and JSCC systems allocate the shared channel by each user's
predicted quality: channel assignment and the number of transmitted semantic symbols from a
semantic spectral efficiency \cite{Yan22}; user association and bandwidth from a message-level
throughput \cite{Xia24}; compression ratios and transmission rates by Bayesian optimization over
Gaussian-process models of quality against SNR and compression ratio \cite{Hou26}; and, in a
non-orthogonal uplink shared with a bit-based user, the communication mode of the semantic user
in each fading state \cite{Mu23}. Multi-user semantic transceivers are also trained end to end for
image retrieval, machine translation and visual question answering \cite{Xie22}. Where the
allocation changes over time, it is learned by deep reinforcement learning: proximal policy
optimization with an attention network assigns resource blocks and the part of each user's
semantic content to send \cite{Wang22}, and a deep deterministic policy gradient agent sets each
user's compression ratio, transmit power and bandwidth \cite{Zhang23t}. The allocations value
each user's transmission at one SNR and decide it as a whole, so none of them can value a slot
for an image whose symbols are spread over slots of different SNRs; the learned policies need
many interactions to train and come with no bound on their distance from the optimum.

\subsection{Restless Bandits and Lagrangian Index Policies}
\label{sec:related-rmab}

Whittle relaxed the constraint on the number of projects operated at each instant to a
constraint on its expected value. The Lagrange multiplier of the relaxed constraint defines an
index, and Whittle conjectured that operating the projects of largest index is nearly optimal
when there are many of them \cite{Whittle88}. Weber and Weiss proved the conjecture when the
fluid approximation of the index policy has a globally stable equilibrium point, and gave an
example in which it fails \cite{WeberWeiss90}. More generally, dynamic programs whose subproblems
are coupled only through constraints on their actions admit Lagrangian and linear-programming
relaxations whose value functions separate by subproblem \cite{Adelman08}. Relaxing the limit on
the number of items selected at a time yields both an upper bound on the optimal value and the
Lagrangian index policy, and under mild conditions both are asymptotically optimal with many
items, with tie-breaking playing an essential role \cite{Brown20}. A linear-program-based index
policy is asymptotically optimal over finite and infinite horizons, at a square-root rate in
general and exponentially fast under further conditions \cite{Gast24}. For deadline scheduling,
Hao, Xu and Tong define the Lagrangian priority value of a job as the greatest tax under which
it is optimal to process it, and prove that serving jobs in that order is asymptotically optimal
as the system scales \cite{Hao22}.

The same tools are established in wireless scheduling. For dynamic multichannel access, the
Whittle index has a closed form, and for non-identical channels a Lagrangian relaxation gives an
upper bound on performance \cite{LiuZhao10}. For the age of information of clients behind
unreliable channels, Whittle's index policy performs near optimally \cite{Kadota18}. For remote
inference, the joint choice of feature lengths and transmissions forms a restless multi-armed,
multi-action bandit \cite{Shisher23}. Closest to ours, McDowell et al. schedule a learned semantic
code over independent Rayleigh-fading slots: a Lagrangian relaxation with one price splits the
problem into per-sensor programs whose values are learned by reinforcement learning, the reward is
a task value computed from the decoder's posterior, and a fluid upper bound is stated but not
plotted \cite{McDowell26}. Each slot, however, carries a fresh codeword of the newest observation,
so nothing accumulates across slots toward a deadline, and the channel state is known exactly.

Restless bandits thus already schedule layered video \cite{Hosseini19} and learned semantic codes
\cite{McDowell26}, which supports casting the deadline scheduling of learned codes as one. CLASP
imports this machinery: the Lagrangian relaxation with one price per slot, its upper bound, and the
Lagrangian priority policy \cite{Brown20,Hao22}. What is new is the arm and the reward:
each user's phases accumulate over slots of different SNRs toward a decoded quality at the
deadline, the link abstraction reduces each user to a two-number state whose program is solved
exactly, and the reward accounts for the measured error of the quality prediction.
